\documentclass[chiro,pdf]{prosper}
\usepackage{url}
\usepackage{sfmath}  % sans serif in math mode.

% Toegegeven, LaTeX is misschien niet de meest laagdrempelige manier
% om slides te maken, maar het laat zich wel heel gemakkelijk
% in versiecontrole gieten.


\title{GroepsAdministratieProgramma}
\author{Introductie, stand van zaken en demo}
\email{informatica@chiro.be}
\institution{
	Chirojeugd Vlaanderen \\
	Kipdorp 30 \\
	2000 Antwerpen
}

\slideCaption{\textit{Introductie, stand van zaken en demo}}
\Logo{\includegraphics[height=15mm]{images/logo.eps}}


\begin{document}

\maketitle

%
% OVERZICHT
%

\overlays{7}
{
\begin{slide}{Overzicht}

\begin{itemstep}
\item Afbakening project
\item Demonstratie
\item Wat verandert er in Kipadmin?
\item Ondersteuning voor de groepen
\item Een aantal cijfers
\item `Hete' problemen
\item Vragen
\end{itemstep}

\end{slide}
}

%
% PROJECTAFBAKENING
%

\overlays{6}
{
\begin{slide}{Projectafbakening}

De functionaliteit van het nieuwe programma wordt teruggebracht tot de essentie:
\begin{itemstep}
\item Beheer van leden voor de groep
\item Aansluiten van leden en leiding
\item Verzekering loonverlies
\item Dubbelpuntabonnementen
\item Bivakaangifte
\item E-mailadressen in en uit maillijsten
\end{itemstep}

\end{slide}
}

%
% DEMO
%

\overlays{12}
{
\begin{slide}{Demo}

\begin{itemize}
\item Webtoepassing: \url{https://gap.chiro.be}
\onlySlide*{2}
{
\begin{itemize}
\item geen installatie nodig
\item meerdere personen kunnen tegelijk werken
\item platformonafhankelijk (Windows, Linux, Mac, \ldots)
\end{itemize}
}
\fromSlide{3}{\item Initi\"ele gegevens}
\fromSlide{4}{\item Personen en ingeschreven leden en leiding}
\fromSlide{5}{\item Automatische jaarovergang en afdelingensysteem}
\fromSlide{6}{\item Personen in- en uitschrijven; probeerperiode}
\fromSlide{7}{\item Lijsten en etiketten via Excel}
\fromSlide{8}{\item Adressensysteem}
\onlySlide*{9}{
\begin{itemize}
\item meerdere adressen per persoon
\item ingebouwde stratenlijst
\end{itemize}
}
\fromSlide{10}{\item Gezinnen beheren}
\onlySlide*{11}{
\begin{itemize}
\item broertjes en zusjes maken
\item een heel gezin tegelijk verhuizen
\end{itemize}
}
\fromSlide{12}{\item Eigen categorie\"en en functies}
\end{itemize}

\end{slide}
}


%
% WAT VERANDERT ER IN KIPADMIN
%

\overlays{7}
{
\begin{slide}{Wat verandert er in Kipadmin?}

\begin{itemize}
\fromSlide{2}{\item Automatische aanmaak van `aansluitingen' voor leden waarvan de probeerperiode voorbij is}
\fromSlide{3}{\item Wijzigingen van naam, adres, geboortedatum, telefoonnummer, \ldots van personen met AD-nummer, komen (quasi) direct door}
\fromSlide{4}{\item Dubbelpuntabonnementen en verzekering loonverlies komen (quasi) direct door}
\fromSlide{5} {\item \'Alle leden in Kipadmin; niet enkel keti's en aspi's}
\fromSlide{6} {\item Aansluitingen en verzekering loonverlies m\'oeten nu via GAP}
\fromSlide{7} {\item Bijwerken van de stratenlijst zal op het secretariaat moeten gebeuren}
\end{itemize}

\end{slide}
}

%
% ONDERSTEUNING
%

\overlays{4}
{
\begin{slide}{Ondersteuning voor de groepen}

\begin{itemstep}
\item Ingebouwde handleiding
\item Chiroforum: \url{http://forum.chiro.be/forum/144}
\item Helpdesk: {helpdesk@chiro.be}
\item Telefoon: 03-231 07 95
\end{itemstep}

\end{slide}
}


%
% EEN AANTAL CIJFERS
%

\overlays{5}
{
\begin{slide}{Cijfers 19 oktober 2010}

\begin{itemstep}
\item Gemiddeld 1728 hits op de webserver per uur
\item Browsers
\begin{itemize}
\item Internet Explorer ($44.5 \%$)
\item Firefox ($35.8 \%$)
\item Chrome ($14.2 \%$)
\item Safari ($4.9 \%$)
\end{itemize}
\item Besturingssystemen
\begin{itemize}
\item Windows ($93.2 \%$)
\item Mac ($5.3 \%$)
\item Linux ($1.5 \%$)
\end{itemize}
\item{683 groepen die jaarovergang uitvoerden}
\item{50571 leden (9223 vorig jaar 19/10)}
\end{itemstep}

\end{slide}
}

%
% TODO
%

\overlays{7}
{
\begin{slide}{`Hete' problemen}
\begin{itemize}
\fromSlide{1}{\item Gewesten en verbonden}
\fromSlide{2}{\item Probleem export naar Excel}
\fromSlide{3}{\item Mails tegen einde probeerperiode}
\fromSlide{4}{\item Inschrijvingen voor bivak}
\fromSlide{5}{\item Interface om stratenlijst bij te werken}
\fromSlide{6}{\item Buitenlandse adressen}
\fromSlide{7}{\item GAV's bijmaken}
\end{itemize}
\end{slide}
}


%
% VRAGEN
%

\begin{slide}{Vragen?}
\end{slide}

\end{document}
